\chapter{Engineering Interpretation and Trade Study}
\label{ch:engineering-interpretation}
\textbf{Chapter status: PENDING RESULTS; interpretation questions frozen.}

This chapter is mandatory. It does not repeat the result tables; it explains why the observed
behavior occurred and what engineering decisions follow.

\section{Numerical Accuracy versus Word Length}
Interpret the relationship among quantization step, covariance error, structural preservation,
and downstream DOA consistency.

\section{Accuracy versus FPGA Resources}
Compare numerical benefit against DSP, LUT, FF, BRAM, and routing cost. The preferred
format shall be defended, not merely reported.

\section{Latency versus Throughput}
Latency and initiation interval/throughput shall be discussed separately. A design with lower
latency is not automatically the design with higher sustained frame throughput.

\section{Pipeline and Unroll Interpretation}
For each pragma or architecture change, explain the causal chain from dependency/memory
access to II/latency and resource usage.

\section{Timing Closure Interpretation}
Report requested clock, achieved timing, WNS, critical path, evaluated mitigation, and final
decision. A missed target is an engineering result, not a reason to hide the run.

\section{Integration Interpretation}
Separate numerical discrepancies from packing/order/interface faults using the bit-accurate
reference hierarchy.

\section{End-to-End Engineering Question}
The final trade study shall answer:
\begin{equation}
\boxed{\text{What FPGA implementation cost is required to preserve useful DOA fidelity?}}
\end{equation}
The answer must be bounded by the actual EE 6800 baseline and must not imply full radar or
full MUSIC hardware acceleration.

\section{E2E-G2 --- Engineering Interpretation Gate}
\textbf{Exit criterion.} Major numerical, architecture, timing, resource, and integration
results have causal explanations and explicit design decisions.\\
\textbf{Status.} PENDING.
